\hypertarget{text-measurement-must-preserve-the-causal-question}{%
\subsection{Text measurement must preserve the causal
question}\label{text-measurement-must-preserve-the-causal-question}}

\label{exec13:census-measurement}

A servicing analyst may ask an assistant to extract financial distress
from call notes and then adjust for the extracted score when comparing
repayment plans. This can be useful work. It also combines three
different claims: that the score measures distress, that the note was
available before the plan decision, and that conditioning on the score
removes the relevant confounding. An accurate summary of the note
establishes none of the latter two claims. A later note may describe a
reaction to the plan; a note written only after a complaint may select
on both treatment and outcome. The proposed formulation interface must
attach a time and a measurement interpretation to each extracted
variable before adding it to an adjustment set.

Keith, Jensen and O'Connor review text used to measure otherwise
unobserved confounders. Their discussion explicitly distinguishes
confounders, colliders and mediators and asks whether one document can
encode several of these at once. It also identifies error in measured
confounders as an unresolved problem. A richer language representation
can improve measurement, but the review provides no guarantee that a
representation eliminates confounding. The appropriate comparison
therefore includes simple dictionaries, supervised extractors,
human-coded measurements and a language model, with the same permitted
document dates and the same effect target.
\citep{exec13_census_text_review_arxiv}

\hypertarget{a-reproducible-example-of-residual-confounding}{%
\subsubsection{A reproducible example of residual
confounding}\label{a-reproducible-example-of-residual-confounding}}

Consider the following project-generated population model. The
intervention intensity \(A\) and later loss \(Y\) are standardized
continuous quantities; these are illustrative units, not calibrated bank
measurements. True predecision distress is \(U\), while \(W\) is an
imperfect score extracted from a predecision note:

\[
 A=U+\varepsilon_A,\qquad
 W=U+\varepsilon_W,\qquad
 Y=\tau A+U+\varepsilon_Y.
\]

All four disturbances are mutually independent and centered Gaussian.
The variances of \(U,\varepsilon_A,\varepsilon_Y\) are one, while
\(\operatorname{Var}(\varepsilon_W)=s\geq0\). The true effect of
increasing \(A\) by one is \(\tau\). Conditional on \(W\), the residual
variance of \(U\) is \(v=s/(1+s)\). Residualizing both \(A\) and \(Y\)
on \(W\) therefore gives the population regression coefficient

\[
 \beta_{A\mid W}
 =\tau+\frac{v}{1+v}
 =\tau+\frac{s}{1+2s}.
\]

The second term persists with an infinite sample. It is not estimation
noise that a larger language model, an accurate propensity fit to \(W\),
or a narrower regression interval automatically removes. With the
stipulated \(\tau=-1/5\), exact covariance calculations give the
following values.

\begin{longtable}[]{@{}lll@{}}
\toprule
Text-score noise variance \(s\) & True effect & Adjusted population
coefficient\tabularnewline
\midrule
\endhead
\(0\) & \(-1/5\) & \(-1/5\)\tabularnewline
\(1/4\) & \(-1/5\) & \(-1/30\)\tabularnewline
\(1\) & \(-1/5\) & \(2/15\)\tabularnewline
\(4\) & \(-1/5\) & \(11/45\)\tabularnewline
\bottomrule
\end{longtable}

The score with \(s=1\) leaves enough confounding to reverse the sign.
This example is independently reproduced by solving the population
covariance normal equations and by residual variance decomposition. It
demonstrates a specific failure mechanism under a stated law, not a
claim that all text adjustment fails. A successful extraction study must
measure the relevant error and evaluate how it affects the causal
target. If a proxy-correction method is proposed, its additional
measurement assumptions and available validation information must be
stated rather than inferred from predictive accuracy.

\hypertarget{freeze-the-measurement-rule-before-evaluating-its-effects}{%
\subsubsection{Freeze the measurement rule before evaluating its
effects}\label{freeze-the-measurement-rule-before-evaluating-its-effects}}

Egami and colleagues make the measurement rule itself explicit. Let
\(g\) map a document into a category or vector. For a textual outcome,
\(E[g(Y(1))-g(Y(0))]\) is a target defined by that particular rule.
Searching over rules using the same treatment assignments and documents
later used for effect estimation can induce dependence and select
apparent effects. Their proposed separation of codebook discovery from
effect estimation addresses that problem, while their text-as-treatment
discussion imposes a further sufficiency restriction on the
representation. \citep{exec13_census_texts}

For a language-based measurement pipeline, the rule includes the model
version, prompt, retrieval corpus, preprocessing, extraction schema and
response aggregation. Freezing only a classifier filename while changing
the prompt after reading the effect estimates does not freeze \(g\). The
prototype should develop the rule on a designated development set,
record its identity, and evaluate it on held-out units. If several calls
concern one customer or conversation, those calls should remain together
in the split. This is a proposed implementation of the separation
principle, not a claim that sample splitting alone ensures valid
inference for dependent histories.

The held-out review asks whether the labels retain their intended
meaning, whether extraction errors differ by treatment and outcome, and
whether missing documents have a defensible observation model. It must
preserve negative findings. Rewriting the measurement rule to rescue a
disappointing effect requires a new evaluation set or an explicitly
exploratory analysis. Neither a fixed rule nor a held-out set supplies
unmeasured assignment information or repairs a post-treatment adjustment
error.

\hypertarget{formulate-competing-explanations-before-selecting-a-diagram}{%
\subsubsection{Formulate competing explanations before selecting a
diagram}\label{formulate-competing-explanations-before-selecting-a-diagram}}

Geraldo Bastías proposes building diagrams from the assumptions actually
used in a study, preserving several models where the assumptions do not
determine one. His workflow checks identification within each model,
derives observable implications, and distinguishes missing relationships
from missing variables. This is a methodological contribution to diagram
review, not an implemented LLM discovery algorithm. It supports an
assistant that makes assumptions inspectable rather than treating a
familiar design label as sufficient. \citep{exec13_census_credible}

For the servicing example, a short initial brief can leave open whether
distress caused the plan assignment, whether a later complaint caused a
note to be recorded, and whether the proposed outcome is arrears, loss
or customer welfare. The assistant should construct separate compatible
diagrams and show how each changes adjustment or identification. It can
check timing contradictions and derive conditional-independence
implications. Passing those checks does not establish that an omitted
common cause is absent. Where the diagrams imply different effects, the
useful output is the unresolved choice and the evidence that could
distinguish it.

These requirements permit immediate public and generated evaluation. A
public document corpus can supply realistic language, while a declared
generator supplies assignment, document timing, measurement error and
outcome truth. Simple extraction and fixed numerical workflows provide
baselines. The evaluation must vary mechanisms that are difficult for
the chosen extractor, including selective note availability and notes
written after treatment. Keith and colleagues specifically caution that
authors can choose semi-synthetic confounders their own method can
easily recover. Freezing adverse mechanisms and separating their
construction from final method selection addresses that concern more
directly than adding familiar domain vocabulary to a favorable
simulation. \citep{exec13_census_text_review_arxiv}

\label{exec13:census-measurement-end}
